\documentclass[10pt,a4paper]{amsart}
\usepackage[margin=19mm]{geometry}
\usepackage{amsmath,amssymb,mathtools}
\usepackage[scr=rsfs,cal=euler]{mathalfa}
\usepackage{xfrac}
\usepackage[unicode,hidelinks,bookmarks=false]{hyperref}
\newcommand{\ie}{i.e.\ }
\newcommand{\cf}{cf.\ }
\newcommand{\resp}{respectively\ }
\newcommand{\Subsection}{\subsection}
\providecommand{\nobreakdash}{\nobreak\hbox{-}}
\setlength{\emergencystretch}{2em}
\multlinegap0pt
\usepackage{bm}
\usepackage{bbm}
\usepackage{dsfont}
\usepackage{xspace}
\usepackage{cancel}
\usepackage{mathtools}
\usepackage{amsthm}
\makeatletter
\@ifundefined{theorem}{\newtheorem{theorem}{Theorem}}{}
\@ifundefined{lemma}{\newtheorem{lemma}[theorem]{Lemma}}{}
\makeatother
\title[The Riemannian Geometry Associated to Gradient Flows of Line]{The Riemannian Geometry Associated to Gradient Flows of Linear Convolutional Networks}
\author{El Mehdi Achour, Kathl\'en Kohn, Holger Rauhut}
\date{Source: arXiv:2507.06367v2 (CC BY 4.0)}
\begin{document}
\maketitle

\noindent\emph{Source and licence.} The Riemannian Geometry Associated to Gradient Flows of Linear Convolutional Networks. Version arXiv:2507.06367v2 (CC BY 4.0), distributed under the Creative Commons Attribution 4.0 International licence, as recorded in the arXiv licence metadata for this exact version and shown on the abstract page https://arxiv.org/abs/2507.06367v2. Full rights notice: licenses/arxiv-2507.06367-CC-BY-4.0.txt.\par\smallskip

\noindent\textit{Editorial scope.} Counted unit: the Lemma whose source label is \texttt{lem:technical} in the article (source0.tex lines 1367--1371, source label \texttt{lem:technical}), delivered together with the article's own complete original proof of it (source0.tex lines 1374--1394, one \texttt{proof} environment of 21 lines). No mathematics is written, repaired or re-derived: statement and proof are the article's own text line for line. Required context retained uncounted: none. Every definition, display and label of those lines is retained unchanged. Omitted from this package and not redistributed: 1--1366, 1372--1373, 1395--1584. Nothing inside a retained excerpt is omitted or altered: every retained body is delivered exactly as the article's lines are written, so no omission receipt of any kind accompanies this package. Citations: the retained excerpts carry no citation command at all, so no bibliography entry of the article is transcribed and no reference is redistributed. Reference adaptation: none was needed and none was made; every reference inside the delivered unit resolves inside it, so no unresolved reference or citation is produced. Printed slips kept literal: no printed word, symbol, hypothesis or number of the counted statement or of its proof was altered. Layout and class: the article's own class is \texttt{article} with packages including aistats2026, natbib, graphicx, amsmath, inputenc, babel, fontenc, hyperref, url, booktabs, amsfonts, nicefrac, microtype, xcolor; the wrapper is the lane's pinned \texttt{amsart} editorial wrapper at 10pt on A4, which restates the theorem-like environments the retained text needs and adds the article's own macro definitions that the retained text needs (none), with extra wrapper packages (bm, bbm, dsfont, xspace, cancel, mathtools). The delivered numbering is the unit's own. Assets: no retained span contains a figure environment, a graphics-inclusion command, a PDF-page inclusion, a table or a TikZ picture; no image file is copied into this package, the package declares no asset dependency and no image file is redistributed with it. Licence: version arXiv:2507.06367v2 of this article is distributed under the Creative Commons Attribution 4.0 International licence, as recorded in the arXiv licence metadata for this exact version and shown on the abstract page https://arxiv.org/abs/2507.06367v2. A full rights notice is delivered at licenses/arxiv-2507.06367-CC-BY-4.0.txt. Characters: every retained excerpt is pure ASCII, so the delivered engine, XeLaTeX with its default Latin Modern text font, reports no missing character.
\begin{lemma}
\label{lem:technical}
    Let $\theta = (w_1, \ldots, w_H)$ with $w_i \in \mathbb{R}^{k_i} \setminus \{ 0 \}$, and let $N \geq k$.
    For almost all $N$-tuples of data $\mathcal{D}$, the vector $Au$ is not orthogonal to the image of the derivative $d_\theta \mu$.
\end{lemma}

\begin{proof}
    Let us first consider the case that $\mu(\theta)$ is a smooth point of the neuromanifold $\mathcal{M}$.
    Then, the image of $d_\theta \mu$ is the tangent space of $\mathcal{M}$ at $\mu(\theta)$.
    A vector $v=(v_0, \ldots, v_{k-1})$ is orthogonal to this tangent space if and only if the hyperplane 
    $P_v := \{ (h_0, \ldots, h_{k-1}) \mid \sum_i v_i h_i=0 \}$ contains the tangent space.
    Due to the genericity of the data $\mathcal{D}$, the vector $Au$ and the hyperplane $P_{Au}$ are generic. 
    Since the dual variety of the Zariski closure $\mathbb{P}\overline{\mathcal{M}}$ of the projective neuromanifold $\mathbb{P}\mathcal{M}$ is a proper subvariety of $(\mathbb{P}^{k-1})^\ast$, the generic hyperplane $P_{Au}$, considered as a point in the dual space $(\mathbb{P}^{k-1})^\ast$, is not contained in the dual variety of $\mathbb{P}\overline{\mathcal{M}}$.
    Hence, the hyperplane $P_{Au}$ is not tangent to $\mathbb{P}\overline{\mathcal{M}}$, and so the vector $Au$ is not orthogonal 
    to $T_{\mu(\theta)} \mathcal{M} = \mathrm{im}(d_\theta \mu)$.
    
    Now, we consider  arbitrary $\theta$ where all layer filters are nonzero.
    We denote by $r$ the rank of the derivative $d_\theta \mu$.
    Since $d_\theta \mu (\theta) = H \cdot  \mu(\theta) \neq 0$, we have that $r \geq 1$.
    We further denote by $X_r$ the set of all filter tuples $\theta'$ such that the rank of $d_{\theta'}\mu$ is $r$, and by $Y_r$ the Zariski closure of $\mu(X_r)$.
    Note that $Y_{\dim \mathcal{M}} = \overline{\mathcal{M}}$.
    If $\mu(\theta)$ is a smooth point of $Y_r$, then $\mathrm{im}(d_\theta \mu)$ is the tangent space $T_{\mu(\theta)} Y_r$.
    Otherwise, since $\mathrm{im}(d_\theta \mu)$ has dimension $r$, this linear space is contained in the Zariski closure of the set of all tangent spaces at smooth points of $Y_r$.
    Therefore, the dual variety of $\mathbb{P} Y_r$ contains all hyperplanes that contain the $r$-dimensional $\mathrm{im}(d_\theta \mu)$.
    Since, as above, the hyperplane $P_{Au}$ is a generic point in $(\mathbb{P}^{k-1})^\ast$, it is not contained in the dual variety of $\mathbb{P} Y_r$.
    Thus, it does not contain $\mathrm{im}(d_\theta \mu)$, and so $Au$ is not orthogonal to $\mathrm{im}(d_\theta \mu)$.
\end{proof}

\end{document}
